\section{Acesso aos Dados}
\label{sec:acesso}

\subsection{Pela API}

A API abre um \emph{pool} de até dez conexões com o driver \campo{pg}, com
tempo limite de 5 segundos para conectar e de 10 segundos por
consulta. Cada repositório em \campo{api/adapters/db} escreve o próprio SQL,
com parâmetros (\campo{\$1}, \campo{\$2}) em vez de concatenação de texto, o
que impede injeção de SQL.

Operações com mais de um passo rodam numa transação, pela função
\campo{inTransaction}: ou tudo é gravado, ou nada. É o caso de salvar uma música,
que grava a linha em \campo{musics} e regrava os artistas, os gêneros e a faixa
no álbum, e de criar um convite do plano familiar.

O convite mostra por que só a transação não basta. O limite de quatro pessoas
conta os membros e os convites pendentes, e duas requisições simultâneas
poderiam contar a mesma vaga livre e criar um convite cada uma. Para evitar
isso, a transação começa com um \emph{lock} consultivo do Postgres
(\campo{pg\_advisory\_xact\_lock}) com o \campo{id} do plano: o segundo
convite espera o primeiro terminar e já enxerga a vaga ocupada.

As colunas \campo{date}, como a data de nascimento, voltam como texto
(\campo{AAAA-MM-DD}). O conversor padrão do driver as transformaria em
meia-noite no fuso do servidor, e a data poderia mudar de dia. Isso afetaria
a regra de idade do conteúdo explícito.

\subsection{Seed}

O seed, em \campo{api/prisma/seed/}, popula o banco com o client do Prisma.
Cada seeder tem um ambiente:

\begin{itemize}
  \item \textbf{ALL}: dados reais do produto, como os tipos de plano e os
    gêneros. Rodam no banco local e em produção, e nunca apagam nada.
  \item \textbf{LOCAL}: dados fictícios para desenvolver e testar (usuários,
    catálogo, playlists e interações). Só rodam num banco em
    \campo{localhost}.
\end{itemize}

O seed roda numa transação e usa uma semente fixa, então gera os mesmos dados
e os mesmos \campo{id}s em todas as máquinas. Isso torna os testes manuais
reproduzíveis entre os integrantes.

\subsection{Ambientes}

Localmente, o Postgres roda num contêiner Docker, acessível só pela própria
máquina. Em produção, o banco roda no servidor da equipe e não é exposto à
internet: o acesso passa por um túnel da Cloudflare protegido pelo Cloudflare
Access, liberado só para os emails da equipe. Os usuários do Postgres da
equipe podem ler e alterar dados, mas não mudar o esquema, que é papel das
migrations, nem ler as tabelas com tokens e hashes de senha.
